%=============================================================
% Capitolo 5 — Condizionamento dei Sistemi Lineari
%=============================================================
\chapter{Condizionamento dei Sistemi Lineari}
\label{ch:condizionamento}

\section{Il problema della perturbazione}

Nella pratica i dati di un sistema lineare $Ax = b$ non sono mai noti con precisione
infinita: sono affetti da errori di misura, di arrotondamento, o di approssimazione.
È fondamentale capire \emph{quanto} la soluzione $x$ cambia al variare (anche
leggermente) dei dati $A$ e $b$.

%------------------------------------------------------------
\section{Perturbazione solo sul termine noto}
\label{sec:pert_b}
%------------------------------------------------------------

Consideriamo il sistema perturbato
\begin{equation}
  A(x + \delta x) = b + \delta b,
  \label{eq:pert_b}
\end{equation}
dove $\delta b$ è una piccola perturbazione del termine noto e $\delta x$ è la
conseguente perturbazione della soluzione. Sottraendo $Ax = b$:
\begin{equation}
  A\,\delta x = \delta b \implies \delta x = A^{-1}\delta b.
\end{equation}

Prendendo le norme:
\begin{equation}
  \|\delta x\| \leq \|A^{-1}\|\,\|\delta b\|.
  \label{eq:bound_deltax}
\end{equation}
Poiché $\|b\| = \|Ax\| \leq \|A\|\,\|x\|$, si ottiene $\frac{1}{\|x\|} \leq
\frac{\|A\|}{\|b\|}$. Combinando:

\begin{equation}
  \frac{\|\delta x\|}{\|x\|} \leq \|A\|\,\|A^{-1}\|\,\frac{\|\delta b\|}{\|b\|}
  = \kappa(A)\,\frac{\|\delta b\|}{\|b\|}.
  \label{eq:cond_b}
\end{equation}

%------------------------------------------------------------
\section{Perturbazione sulla matrice e sul termine noto}
\label{sec:pert_Ab}
%------------------------------------------------------------

Consideriamo ora il caso più generale in cui sia $A$ che $b$ sono perturbati:
\begin{equation}
  (A + \delta A)(x + \delta x) = b + \delta b.
  \label{eq:pert_Ab}
\end{equation}

Espandendo e trascurando i termini del secondo ordine ($\delta A\,\delta x \approx 0$):
\begin{equation}
  A\,\delta x \approx \delta b - \delta A\, x,
\end{equation}
da cui $\delta x = A^{-1}(\delta b - \delta A\, x)$, e prendendo le norme:

\begin{theorem}[Limitazione dell'errore relativo — caso generale]
\label{thm:cond_gen}
Se $\|\delta A\| < \|A^{-1}\|^{-1}$ (perturbazione piccola), allora:
\begin{equation}
  \frac{\|\delta x\|}{\|x + \delta x\|} \leq
  \frac{\kappa(A)}{1 - \kappa(A)\,\frac{\|\delta A\|}{\|A\|}}
  \left(\frac{\|\delta A\|}{\|A\|} + \frac{\|\delta b\|}{\|b\|}\right).
  \label{eq:cond_gen}
\end{equation}
\end{theorem}

\begin{remark}
Il \cref{thm:cond_gen} mostra che l'errore relativo sulla soluzione è amplificato dal
fattore $\kappa(A)$: se $\kappa(A) = 10^k$, si perdono fino a $k$ cifre significative
nella soluzione rispetto alla precisione dei dati.
\end{remark}

%------------------------------------------------------------
\section{Interpretazione pratica del condizionamento}
\label{sec:interpretazione}
%------------------------------------------------------------

In aritmetica in virgola mobile con epsilon di macchina $\eps_M$, le perturbazioni
sui dati sono dell'ordine di $\eps_M$. La \cref{eq:cond_gen} implica che l'errore
relativo sulla soluzione è approssimativamente:
\begin{equation}
  \frac{\|\delta x\|}{\|x\|} \approx \kappa(A)\,\eps_M.
  \label{eq:errore_pratico}
\end{equation}

\begin{example}[Sistema mal condizionato]
La \textbf{matrice di Hilbert} di ordine $n$, con elementi $h_{ij} = 1/(i+j-1)$,
è notoriamente mal condizionata:
\[
  H_5 = \begin{pmatrix}
    1 & \frac{1}{2} & \frac{1}{3} & \frac{1}{4} & \frac{1}{5} \\[4pt]
    \frac{1}{2} & \frac{1}{3} & \frac{1}{4} & \frac{1}{5} & \frac{1}{6} \\[4pt]
    \frac{1}{3} & \frac{1}{4} & \frac{1}{5} & \frac{1}{6} & \frac{1}{7} \\[4pt]
    \frac{1}{4} & \frac{1}{5} & \frac{1}{6} & \frac{1}{7} & \frac{1}{8} \\[4pt]
    \frac{1}{5} & \frac{1}{6} & \frac{1}{7} & \frac{1}{8} & \frac{1}{9}
  \end{pmatrix}, \quad \kappa_2(H_5) \approx 4.77 \times 10^5.
\]
Con doppia precisione ($\eps_M \approx 10^{-16}$) l'errore relativo sulla soluzione
è dell'ordine di $4.77 \times 10^5 \times 10^{-16} \approx 10^{-11}$: si perdono
circa 5 cifre significative.
\end{example}

%------------------------------------------------------------
\section{Stima del numero di condizionamento}
\label{sec:stima_cond}
%------------------------------------------------------------

Calcolare $\kappa(A) = \|A\|\,\|A^{-1}\|$ esattamente richiede di conoscere $A^{-1}$,
il che è costoso ($\mathcal{O}(n^3)$). In pratica:
\begin{itemize}
  \item Per la norma 2: $\kappa_2(A) = \sigma_{\max}(A)/\sigma_{\min}(A)$
        (rapporto tra valore singolare massimo e minimo), calcolabile con la SVD.
  \item Metodi di stima a basso costo (e.g.\ \texttt{rcond} in MATLAB/NumPy)
        forniscono una stima di $1/\kappa(A)$ in $\mathcal{O}(n^2)$ una volta
        calcolata la fattorizzazione LU.
\end{itemize}

%------------------------------------------------------------
\section{Residuo e soluzione calcolata}
\label{sec:residuo}
%------------------------------------------------------------

Sia $\hat{x}$ la soluzione calcolata numericamente di $Ax = b$. Il
\textbf{residuo} è:
\begin{equation}
  r = b - A\hat{x}.
  \label{eq:residuo}
\end{equation}

Un residuo piccolo \textbf{non} implica necessariamente un errore piccolo. Il legame
è:
\begin{equation}
  \frac{\|x - \hat{x}\|}{\|x\|} \leq \kappa(A)\,\frac{\|r\|}{\|b\|}.
  \label{eq:residuo_errore}
\end{equation}

\begin{remark}[Backward e forward error]
\begin{itemize}
  \item Il \emph{backward error} (errore all'indietro) è $\|r\|/\|b\|$: misura di
        quanto $\hat{x}$ risolve \emph{esattamente} un sistema perturbato.
  \item Il \emph{forward error} (errore in avanti) è $\|x - \hat{x}\|/\|x\|$:
        misura la distanza dalla soluzione vera.
\end{itemize}
Un metodo si dice \textbf{backward-stabile} se il backward error è dell'ordine di
$\eps_M$.
\end{remark}

%------------------------------------------------------------
\section{Riepilogo}
\label{sec:riepilogo_cond}
%------------------------------------------------------------

\begin{table}[htbp]
  \centering
  \caption{Effetto del numero di condizionamento sulla soluzione (doppia precisione,
           $\eps_M \approx 10^{-16}$).}
  \label{tab:condizionamento}
  \begin{tabular}{ccc}
    \toprule
    $\kappa(A)$ & Cifre perse & Qualità \\
    \midrule
    $\sim 1$       & $\sim 0$  & Eccellente \\
    $\sim 10^3$    & $\sim 3$  & Buona \\
    $\sim 10^8$    & $\sim 8$  & Accettabile \\
    $\sim 10^{12}$ & $\sim 12$ & Scarsa \\
    $\sim 10^{16}$ & $\sim 16$ & Soluzione inattendibile \\
    \bottomrule
  \end{tabular}
\end{table}

In sintesi: prima di risolvere un sistema lineare è \emph{sempre} buona pratica
stimare $\kappa(A)$ per capire quante cifre significative ci si può aspettare nella
soluzione.
